\chapter{Le réseau de neurones : précis, mais pas un QUBO}\label{ch:mlp}
\iterbanner{V3}

\section{Pourquoi l'étudier}
Les trois modèles précédents sont quadratiques en bits, parce que le solveur quantique l'exige. Un réseau de neurones ne l'est pas. Le mettre dans la même boucle répond à une question simple : \textbf{combien coûte la contrainte « le modèle doit être un QUBO »} ? Si un modèle libre fait nettement mieux, c'est autant de précision que la formulation quantique oblige à abandonner.

\section{Le modèle}
\begin{definition}[ (perceptron à deux couches cachées)]
Pour une entrée $\mathbf z\in\R^d$ (les variables normalisées, comme le modèle de la V2) :
\[ \mathbf a_1=\tanh(W_1^\top\mathbf z+\mathbf c_1),\qquad\mathbf a_2=\tanh(W_2^\top\mathbf a_1+\mathbf c_2),\qquad\hat y=\mathbf w_3\cdot\mathbf a_2+c_3, \]
avec 32 neurones par couche cachée. La fonction $\tanh$ est appliquée composante par composante.
\end{definition}
Pour $d=4$, cela fait $4\times32+32+32\times32+32+32+1=1\,249$ paramètres, presque cent fois plus que le modèle quadratique en $z$ (15). Pour $d=10$ (problème mis à l'échelle), 1\,441.

\paragraph{L'apprentissage.} Comme pour la machine de factorisation : erreur quadratique moyenne plus pénalité $L_2$ ($\lambda=10^{-3}$), sur la sortie centrée réduite, minimisée par Adam (1\,000 à 1\,500 itérations). Le gradient se calcule couche par couche, de la sortie vers l'entrée (rétropropagation) : avec $\mathbf r=\frac2N(\hat{\mathbf y}-\mathbf y)$,
\[ \frac{\partial\mathcal L}{\partial\mathbf w_3}=A_2^\top\mathbf r,\qquad\boldsymbol\delta_2=(\mathbf r\,\mathbf w_3^\top)\odot(1-A_2^2),\qquad\frac{\partial\mathcal L}{\partial W_2}=A_1^\top\boldsymbol\delta_2,\qquad\boldsymbol\delta_1=(\boldsymbol\delta_2W_2^\top)\odot(1-A_1^2),\ \dots \]
où $1-\tanh^2$ est la dérivée de $\tanh$ et $\odot$ le produit composante par composante.

\paragraph{Un ensemble, pour avoir une incertitude.} Le code entraîne 3 à 5 réseaux à partir de tirages initiaux différents. La prédiction est leur moyenne ; leur écart-type donne une estimation, grossière, de l'incertitude. Là où les données sont nombreuses, les réseaux sont d'accord ; ailleurs, ils divergent.

\section{Pourquoi ce n'est pas un QUBO}
Une fonction de bits s'écrit d'une seule façon comme un polynôme multilinéaire (chaque bit au plus à la puissance 1). Le modèle en $z$ et les modèles en bits n'ont que des termes de degré 0, 1 et 2. Le réseau, lui, compose des $\tanh$ : son développement contient des termes de \emph{tous} les degrés.

\paragraph{Exemple : une seule variable.} Prenons $\hat y=\tanh(3z_m)$ pour la cambrure (8 niveaux, 3 bits). Les 8 valeurs, de $z=-1$ à $z=1$, sont
\[ -0{,}995;\ -0{,}973;\ -0{,}858;\ -0{,}404;\ 0{,}404;\ 0{,}858;\ 0{,}973;\ 0{,}995. \]
Une fonction de 3 bits a 8 coefficients multilinéaires : constante, 3 linéaires, 3 couplages et \emph{un} terme cubique $b_0b_1b_2$. On le trouve en combinant les 8 valeurs avec des signes alternés (exercice~\ref{ex:cubique}) : il ne vaut pas 0 ici. Ce terme n'a pas de place dans un QUBO ; il faudrait un qubit auxiliaire (réduction de Rosenberg) par terme de ce type.

\section{Choisir sans QUBO : la recherche directe}
Sans QUBO, ni l'énumération des chaînes de bits, ni le recuit, ni QAOA ne s'appliquent tels quels. On minimise directement la prédiction sur la grille des profils :
\begin{itemize}
\item \textbf{problème V2} (1\,024 profils) : on évalue le réseau sur toute la grille, ce qui prend quelques millisecondes ;
\item \textbf{problème mis à l'échelle} ($8^{10}\approx10^9$ profils) : on tire 20\,000 profils au hasard, puis, depuis les 20 meilleurs, une \textbf{recherche locale} : on essaie de changer le niveau d'une seule variable (10 variables $\times$ 7 autres niveaux $=70$ voisins) et l'on garde le meilleur voisin, jusqu'à ce qu'aucun ne fasse mieux.
\end{itemize}
Les contraintes sont traitées comme dans le QUBO, sans les défauts de l'encodage : un réseau de plus prédit $C_{L,\max}$, et l'on ajoute $75\,\max(0,\ 1{,}2-\widehat{C_{L,\max}})^2$ (la pénalité idéale, sans les dents de scie des variables d'écart) ; l'épaisseur interdite coûte 3.

\begin{codebox}[ : \code{ia\_qubo/boucle.py}, fonction \code{recherche\_directe}]
\begin{lstlisting}
I = np.column_stack([rng.integers(0, L, n_alea) for L in levels])   # 20 000 profils au hasard
s = score(I)
for x in [I[j] for j in np.argsort(s)[:n_depart]]:                 # 20 departs
    while amelio:                                                   # recherche locale
        voisins = [x avec la variable i au niveau v, pour tous i et v != x[i]]
        j = argmin(score(voisins))
        if score(voisins[j]) < f(x): x = voisins[j]
\end{lstlisting}
\end{codebox}

\begin{projet}
La recherche directe n'a rien de quantique, et rien ne garantit qu'elle trouve le minimum du réseau sur $10^9$ profils. Mais elle est rapide, et l'on n'a pas besoin du minimum exact : on cherche un \emph{bon candidat}, que la simulation départagera. C'est exactement l'argument que l'on oppose à QAOA dans le cours de calcul quantique : la précision du choix compte moins que la qualité du modèle.
\end{projet}

\begin{resultat}
\begin{itemize}
\item Un perceptron à deux couches de 32 neurones a environ 1\,250 paramètres : il représente des effets de tous les ordres, mais demande des données.
\item Ce n'est pas un QUBO : son développement en bits contient des termes de degré 3 et plus.
\item On le minimise par énumération (petite grille) ou par tirage aléatoire et recherche locale (grande grille).
\item Il sert de référence : l'écart entre lui et les modèles quadratiques mesure le prix de la formulation QUBO.
\end{itemize}
\end{resultat}

\exercices
\begin{exercice}\label{ex:mlp-param}
Compter les paramètres d'un perceptron à deux couches de $h$ neurones pour $d$ entrées. Application : $d=4$ et $d=10$, $h=32$. Pour quel $h$ le réseau a-t-il autant de paramètres que le modèle quadratique en $z$ à $d=10$ (66) ?
\end{exercice}
\begin{exercice}\label{ex:cubique}
Pour une fonction $f$ de 3 bits, le coefficient du terme $b_0b_1b_2$ de son polynôme multilinéaire vaut $\sum_{\mathbf b}(-1)^{3-|\mathbf b|}f(\mathbf b)$, où $|\mathbf b|$ est le nombre de bits à 1. (a) Le vérifier sur $f=b_0b_1b_2$. (b) Le calculer pour $\hat y=\tanh(3z_m)$ avec les valeurs du cours (attention : la chaîne $b_0b_1b_2$ correspond à l'indice $x=b_0+2b_1+4b_2$). (c) Que vaut-il pour $\hat y=z_m^2$ ?
\end{exercice}

\begin{pistes}
\begin{itemize}
\item \textbf{Distiller le réseau en QUBO} : ajuster une machine de factorisation sur les prédictions du réseau en des milliers de profils tirés au hasard (des données gratuites), puis résoudre ce QUBO. On garderait la précision du réseau là où une forme quadratique peut la représenter.
\item Remplacer l'ensemble par une vraie méthode d'incertitude (abandon aléatoire, réseaux bayésiens) et l'utiliser pour explorer, comme le tirage de Thompson.
\end{itemize}
\end{pistes}
